\documentclass{ximera}

\begin{document}
	\author{Stitz-Zeager}
	\xmtitle{Exercises for Appendix:  Rational Expressions and Equations}{}

\mfpicnumber{1} \opengraphsfile{ExercisesforAppRatExpEqus} % mfpic settings added 


\label{ExercisesforAppRatExpEqus}

\begin{question}
In Exercises \ref{ratsimpfirst} - \ref{ratsimplast}, perform the indicated operations and simplify.

\begin{problem} \label{ratsimpfirst}
$\frac{x^2-9}{x^2} \cdot \frac{3x}{x^2-x-6}$

\begin{solution}
$\frac{3(x+3)}{x(x+2)}$, $x \neq 3$
\end{solution}
\end{problem}

\begin{problem}
$\frac{t^2-2t}{t^2+1} \div (3t^2 - 2t - 8)$

\begin{solution}
 $\frac{t}{(3t+4)(t^2+1)}$, $t \neq 2$
\end{solution}
\end{problem}

\begin{problem}
$\frac{4y-y^2}{2y+1} \div \frac{y^2-16}{2y^2-5y-3}$

\begin{solution}
$-\frac{y(y-3)}{y+4}$, $y \neq -\frac{1}{2}, 3, 4$ 
\end{solution}
\end{problem}

\begin{problem}
$\frac{x}{3x-1} - \frac{1-x}{3x-1} = \frac{\answer{2x-1}}{3x-1}$
\end{problem}

\begin{problem}
$\frac{2}{w-1} - \frac{w^2+1}{w-1}$

\begin{solution}
$-w-1$, $w \neq 1$
\end{solution}
\end{problem}

\begin{problem}
$\frac{2-y}{3y} - \frac{1-y}{3y} + \frac{y^2-1}{3y} = \frac{\answer{y}}{3}$, $y \neq \answer{0}$
\end{problem}

\begin{problem}
$b+ \frac{1}{b-3} - 2$

\begin{solution}
 $\frac{b^2-5b+7}{b-3}$
\end{solution}
\end{problem}

\begin{problem}
$\frac{2x}{x-4} - \frac{1}{2x+1}$

\begin{solution}
$\frac{4x^2+x+4}{(x-4)(2x+1)}$
\end{solution}
\end{problem}

\begin{problem}
$\frac{m^2}{m^2-4} + \frac{1}{2-m} = \frac{\answer{m+1}}{m+2}$, $m \neq \answer{2}$
\end{problem}

\begin{problem}
$\frac{\frac{2}{x} - 2}{x-1} = -\frac{2}{\answer{x}}$, $x \neq \answer{1}$

\end{problem}

\begin{problem}
$\frac{\frac{3}{2-h} - \frac{3}{2}}{h}$

\begin{solution}
$\frac{3}{4-2h}$, $h \neq 0$
\end{solution}
\end{problem}

\begin{problem}
$\frac{\frac{1}{x+h} - \frac{1}{x}}{h}$

\begin{solution}
 $-\frac{1}{x(x+h)}$, $h \neq 0$
\end{solution}
\end{problem}

\begin{problem}
$3w^{-1} - (3w)^{-1} = \frac{8}{\answer{3w}}$

\end{problem}

\begin{problem}
$-2y^{-1}  + 2(3-y)^{-2}$

\begin{solution}
$-\frac{2(y^2-7y+9)}{y(y-3)^2}$
\end{solution}
\end{problem}

\begin{problem}
$3(x-2)^{-1} - 3x(x-2)^{-2} = -\frac{6}{\answer{(x-2)^2}}$

\end{problem}

\begin{problem}
$\frac{t^{-1} + t^{-2}}{t^{-3}}$

\begin{solution}
$t^2+t$, $t \neq 0$
\end{solution}
\end{problem}

\begin{problem}
$\frac{2(3+h)^{-2} - 2(3)^{-2}}{h}$

\begin{solution}
 $-\frac{2(h+6)}{9(h+3)^2}$, $h \neq 0$ 
\end{solution}
\end{problem}

\begin{problem} \label{ratsimplast}
$\frac{(7-x-h)^{-1} - (7-x)^{-1}}{h}$ 

\begin{solution}
 $\frac{1}{(7-x)(7-x-h)}$, $h \neq 0$ 
\end{solution}
\end{problem}

\end{question}

\begin{question}
In Exercises \ref{rateqnfirst} - \ref{rateqnlast}, find all real solutions.  Be sure to check for extraneous solutions.

\begin{problem} \label{rateqnfirst}
$\frac{x}{5x + 4} = 3$

\begin{solution}
$x = -\frac{6}{7}$
\end{solution}
\end{problem}

\begin{problem}
$\frac{3y - 1}{y^{2} + 1} = 1$

 $y = \answer{1}$ (smaller solution)
 
 $y = \answer{2}$ (larger solution)
\end{problem}

\begin{problem}
$\frac{1}{w + 3} + \frac{1}{w - 3} = \frac{w^{2} - 3}{w^{2} - 9}$

\begin{solution}
$w = -1$ 
\end{solution}
\end{problem}

\begin{problem}
$\frac{2x + 17}{x + 1} = x + 5$

$x=\answer{-6}$ (smaller solution)

$x=\answer{2}$ (larger solution)
\end{problem}

\begin{problem}
$\frac{t^{2} - 2t + 1}{t^{3} + t^{2} - 2t} = 1$

\begin{solution}
No real solutions.
\end{solution}
\end{problem}

\begin{problem}
$\frac{-y^{3} + 4y}{y^{2} - 9} = 4y$ 

\begin{solution}
$y = 0, \pm 2\sqrt{2}$ 
\end{solution}
\end{problem}

\begin{problem}
$w + \sqrt{3} = \frac{3w - w^3}{w - \sqrt{3}}$

$w = \answer{-\sqrt{3}}$ (smaller solution.)

$w=\answer{-1}$ (larger solution)
\end{problem}

\begin{problem}
$\frac{2}{x\sqrt{2} - 1}  - 1 = \frac{3}{x \sqrt{2} + 1}$

\begin{solution}
$x = -\frac{3\sqrt{2}}{2}, \sqrt{2}$
\end{solution}
\end{problem}

\begin{problem} \label{rateqnlast}
$\frac{x^2}{(1 + x\sqrt{3})^2} = 3$ 

$x = \answer{-\frac{\sqrt{3}}{2}}$ (smaller solution)

$x = \answer{-\frac{\sqrt{3}}{4}}$ (larger solution)
\end{problem}

\end{question}

\begin{question}
In Exercises \ref{absratfirst} - \ref{absratlast}, use Theorem \ref{absvalequality} along with the techniques in this section to find all real solutions.

\begin{problem} \label{absratfirst}
$\left|\frac{3n}{n-1}  \right| = 3$ 

$n = \answer{\frac{1}{2}}$
\end{problem}

\begin{problem}
$\left| \frac{2x}{x^2-1}\right| = 2$

\begin{solution}
 $x = \frac{1 \pm \sqrt{5}}{2}, \frac{-1 \pm \sqrt{5}}{2}$
\end{solution}
\end{problem}

\begin{problem}  \label{absratlast}
$\left| \frac{2t}{4-t^2}\right| = \left|\frac{2}{t-2}\right|$

$t = \answer{-1}$
\end{problem}

\end{question}

\begin{question}
In Exercises \ref{solveratcalcfirst} - \ref{solveratcalclast}, find all real solutions and use a calculator to approximate your answers, rounded to two decimal places.

\begin{problem}  \label{solveratcalcfirst}
$2.41 = \frac{0.08}{4 \pi R^2}$

\begin{solution}
$R = \pm \sqrt{\frac{0.08}{9.64 \pi}} \approx \pm 0.05$ 
\end{solution}
\end{problem}

\begin{problem}
$\frac{x^2}{(2.31 -x)^2} = 0.04$

\begin{solution}
$x = -\frac{231}{400} \approx -0.58$, $x = \frac{77}{200} \approx 0.38$ 
\end{solution}
\end{problem}

\begin{problem} \label{solveratcalclast}
$1 - \frac{6.75 \times 10^{16}}{c^2} = \frac{1}{4}$

\begin{solution}
$c = \pm \sqrt{\frac{4 \cdot 6.75 \times 10^{16}}{3}} = \pm 3.00 \times 10^{8}$ 

\textbf{NOTE:}  Challenge yourself - try to get the answer without using a calculator!
\end{solution}
\end{problem}

\end{question}

\begin{question}
In Exercises \ref{litrateqnfirst} - \ref{litrateqnlast}, solve the given equation for the indicated variable.

\begin{problem} \label{litrateqnfirst}
Solve for $y$:  $\frac{1-2y}{y+3} = x$ 

\begin{solution}
 $y = \frac{1 - 3x}{x+2}$, $y \neq -3$, $x \neq -2$
\end{solution}
\end{problem}

\begin{problem}
Solve for $y$: $x = 3 - \frac{2}{1-y}$

\begin{solution}
$y = \frac{x-1}{x-3}$, $y \neq 1$, $x \neq 3$
\end{solution}
\end{problem}

\begin{problem} Solve for $T_{2}$: $\frac{V_{1}}{T_{1}} = \frac{V_{2}}{T_{2}}$

\textbf{NOTE:}  Recall that subscripts on variables have no intrinsic mathematical meaning; they're just used to distinguish one variable from another.  In other words, treat quantities like `$V_{1}$' and `$V_{2}$'  as two different variables as you would `$x$' and `$y$.
\end{problem}

\begin{solution}
$T_{2} = \frac{V_{2}}{T_{1}}{V_{1}}$, $T_{1} \neq 0, T_{2} \neq 0, V_{1} \neq 0$
\end{solution}

\begin{problem}
Solve for $t_{0}$:  $\frac{t_{0}}{1-t_{0}t_{1}} = 2$ 

\begin{solution}
$t_{0} = \frac{2}{2t_{1} + 1}$, $t_{1} \neq -\frac{1}{2}$
\end{solution}
\end{problem}

\begin{problem} 
Solve for $x$:  $\frac{1}{x - v_{r}} + \frac{1}{x + v_{r}} = 5$

\begin{solution}
$x = \frac{1 \pm \sqrt{25v_{r}^2+1}}{5}$, $x \neq \pm v_{r}$
\end{solution}
\end{problem}

\begin{problem} \label{litrateqnlast}
Solve for $R$:  $P = \frac{25R}{(R+4)^2}$ 

\begin{solution}
$R= \frac{-(8P-25) \pm \sqrt{(8P-25)^2 - 64P^2}}{2P} = \frac{(25-8P) \pm 5 \sqrt{25-16P}}{2P}$, $P \neq 0$, $R \neq -4$
\end{solution}
\end{problem}

\end{question}

\end{document}

